% ============================================================
% Section 32: 单量子比特的迹距离与Bloch球
% ============================================================
\section{量子信息：单量子比特的迹距离——Bloch 球上的几何距离}

\begin{modelbox}{题目}
在 $\mathbb{C}^2$（单量子比特）空间中，两个量子态的密度矩阵分别为 $\rho$ 和 $\rho'$。

\textbf{(1)} 用 Bloch 矢量表示，证明迹距离
\[
  D(\rho,\rho') = \frac12\operatorname{Tr}|\rho-\rho'|
\]
满足 $D(\rho,\rho') = \frac12|\bm r - \bm r'|$，其中 $\bm r, \bm r' \in \mathbb{R}^3$ 为 Bloch 矢量。

\textbf{(2)} 给出用密度矩阵元表示的等价形式，并说明取值范围和物理意义。
\end{modelbox}

\vspace{0.3em}
\begin{insightbox}{结论预告}
迹距离在 Bloch 球上恰好等于欧氏距离的一半：$D = \frac12|\bm r - \bm r'|$。最大值 $D=1$ 对应正交纯态（对径点），最小值 $D=0$ 对应同一态。
\end{insightbox}


\subsection{迹距离的定义}

\begin{modelbox}{迹距离 (Trace Distance)}
\[
  D(\rho,\rho') \equiv \frac12\operatorname{Tr}|\rho-\rho'|,
  \quad |A| \equiv \sqrt{A^\dagger A}.
\]
由于 $\rho-\rho'$ 厄米，可对角化。设其特征值为 $\lambda_i$，则 $|\rho-\rho'|$ 的特征值为 $|\lambda_i|$，于是：
\[
  D(\rho,\rho') = \frac12\sum_i |\lambda_i|.
\]
\end{modelbox}


\subsection{Bloch 表示}

\subsubsection{单量子比特密度矩阵}

\begin{tipbox}{Bloch 球参数化}
任意 $2\times 2$ 密度矩阵可写成：
\[
  \rho = \frac12(I + \bm r\cdot\bm\sigma), \qquad
  \rho' = \frac12(I + \bm r'\cdot\bm\sigma),
\]
其中 $\bm r, \bm r' \in \mathbb{R}^3$，$|\bm r| \le 1$，$|\bm r'| \le 1$（等号对应纯态），$\bm\sigma=(\sigma_x,\sigma_y,\sigma_z)$ 为 Pauli 矩阵。
\end{tipbox}

\subsubsection{计算 $\rho-\rho'$}

\begin{align*}
  \rho - \rho' &= \frac12(\bm r - \bm r')\cdot\bm\sigma
  \equiv \frac12 \Delta\bm r \cdot\bm\sigma.
\end{align*}

\subsubsection{$\Delta\bm r\cdot\bm\sigma$ 的特征值}

\begin{derivationbox}{Pauli 矩阵的关键性质}
对任意 $\bm v\in\mathbb{R}^3$，矩阵 $\bm v\cdot\bm\sigma = v_x\sigma_x + v_y\sigma_y + v_z\sigma\sigma_z$ 满足：
\[
  (\bm v\cdot\bm\sigma)^2 = |\bm v|^2 I.
\]
这是因为 $\sigma_i\sigma_j = \delta_{ij}I + i\varepsilon_{ijk}\sigma_k$，且交叉项在平方时抵消。

因此 $(\bm v\cdot\bm\sigma)^2 = |\bm v|^2 I$，其特征值必为 $\pm|\bm v|$（各一个，因为迹为零）。
\end{derivationbox}

应用到 $\bm v = \frac12\Delta\bm r$：
\[
  \rho-\rho' = \frac12\Delta\bm r\cdot\bm\sigma
  \quad\Longrightarrow\quad
  \text{特征值为 } \pm\frac12|\Delta\bm r|.
\]

\subsubsection{求 $|\rho-\rho'|$ 的迹}

$|\rho-\rho'|$ 的特征值是 $|{+\frac12|\Delta\bm r|}| = |{-\frac12|\Delta\bm r|}| = \frac12|\Delta\bm r|$（各一个），所以：
\[
  \operatorname{Tr}|\rho-\rho'| = \frac12|\Delta\bm r| + \frac12|\Delta\bm r| = |\Delta\bm r|.
\]

\begin{formulabox}{Bloch 球上的迹距离}
\[
  \boxed{D(\rho,\rho') = \frac12|\bm r - \bm r'|}
\]
\end{formulabox}


\subsection{矩阵元表示}

\subsubsection{Bloch 分量与矩阵元的对应}

设 $\rho = \begin{pmatrix}\rho_{11} & \rho_{12} \\ \rho_{21} & \rho_{22}\end{pmatrix}$，则：
\[
  x = 2\operatorname{Re}\rho_{12}, \quad
  y = 2\operatorname{Im}\rho_{12}, \quad
  z = \rho_{11} - \rho_{22}.
\]
（由 $\rho = \frac12\begin{pmatrix}1+z & x-iy \\ x+iy & 1-z\end{pmatrix}$ 对比得到。）

\subsubsection{用矩阵元表示迹距离}

\begin{align*}
  |\bm r - \bm r'|^2
  &= (x-x')^2 + (y-y')^2 + (z-z')^2 \\
  &= 4|\rho_{12}-\rho'_{12}|^2 + (\rho_{11}-\rho_{22}-\rho'_{11}+\rho'_{22})^2.
\end{align*}

因此：
\begin{formulabox}{矩阵元形式}
\[
  \boxed{D(\rho,\rho') = \sqrt{|\rho_{12}-\rho'_{12}|^2 + \frac14(\rho_{11}-\rho_{22}-\rho'_{11}+\rho'_{22})^2}}
\]
\end{formulabox}


\subsection{物理意义与取值范围}

\subsubsection{几何解释}

\begin{insightbox}{Bloch 球上的几何图像}
\begin{itemize}[nosep]
  \item $D=0$：$\bm r = \bm r'$，两态完全相同
  \item $0<D<1$：两态部分可区分， Bloch 球上距离越远离区分度越高
  \item $D=1$：$\bm r = -\bm r'$ 且 $|\bm r|=|\bm r'|=1$，即\textbf{正交纯态}（Bloch 球对径点），如 $|0\rangle$ 与 $|1\rangle$
\end{itemize}
\end{insightbox}

\subsubsection{最优测量区分概率}

\begin{insightbox}{迹距离的操作意义}
单次测量下成功区分 $\rho$ 和 $\rho'$ 的最大概率为：
\[
  \boxed{P_{\text{success}}^{\max} = \frac12\bigl(1 + D(\rho,\rho')\bigr)}
\]
\begin{itemize}[nosep]
  \item $D=0$：$P_{\max}=1/2$（随机猜测，无法区分）
  \item $D=1$：$P_{\max}=1$（完美区分）
\end{itemize}
迹距离因此是量子态``可区分性''的度量，在量子密码、量子信道容量等领域有核心应用。
\end{insightbox}

\subsubsection{与保真度的关系}

\begin{tipbox}{迹距离 vs 保真度}
对单量子比特，保真度 $F(\rho,\rho')$ 与迹距离满足 Fuchs--van de Graaf 不等式：
\[
  1 - \sqrt{F} \le D \le \sqrt{1 - F}.
\]
在纯态情况下等号成立：$D = \sqrt{1-F}$。对单量子比特纯态 $|\psi\rangle,|\phi\rangle$，$F=|\langle\psi|\phi\rangle|^2=\cos^2(\theta/2)$（$\theta$ 为 Bloch 矢量夹角），则 $D=\sin(\theta/2)=\frac12|\bm r-\bm r'|$，一致。
\end{tipbox}


\subsection{常见错误与考场策略}

\begin{errorbox}{易错点}
\begin{enumerate}
  \item \textbf{绝对值算符 vs 普通绝对值}：$|A|=\sqrt{A^\dagger A}$ 是对 $A$ 取模（谱分解后特征值取绝对值），不是逐元素的绝对值。对厄米矩阵 $H=U\Lambda U^\dagger$，$|H|=U|\Lambda|U^\dagger$。
  \item \textbf{特征值符号}：$\rho-\rho'$ 的特征值是 $\pm\frac12|\Delta\bm r|$（一正一负），但 $|\rho-\rho'|$ 的特征值都是 $\frac12|\Delta\bm r|$（两个正）。迹距离计算的是后者。
  \item \textbf{纯态条件}：$D=1$ 要求两态为正交\textbf{纯态}。若只是任意两个态，最大值仍为 1，但仅当它们为正交纯态时达到。
  \item \textbf{量纲混淆}：Bloch 矢量 $|\bm r|\le 1$ 无量纲，迹距离 $0\le D\le 1$ 也无量纲。
\end{enumerate}
\end{errorbox}

\begin{cheatbox}{单量子比特迹距离：速查卡}
\textbf{Bloch 表示}：$D = \frac12|\bm r - \bm r'|$

\textbf{矩阵元表示}：
\[
  D = \sqrt{|\Delta\rho_{12}|^2 + \frac14(\Delta\rho_{11}-\Delta\rho_{22})^2}
\]
\textbf{最优区分概率}：$P_{\max} = \frac12(1+D)$

\textbf{Pauli 关键恒等式}：$(\bm v\cdot\bm\sigma)^2 = |\bm v|^2 I$

\textbf{特征值速记}：$\bm v\cdot\bm\sigma$ 的特征值永远是 $\pm|\bm v|$（迹为零，行列式为 $-|v|^2$）
\end{cheatbox}
